% SEGMENT OLP-0028-S01
% Part:sets-functions-relations
% Chapter: size-of-sets
% Section: introduction

\documentclass[../../../include/open-logic-section]{subfiles}

\begin{document}

\olfileid{sfr}{siz}{int}

\olsection{அறிமுகம்}

1870களில் ஜியோர்க் காண்டோர் கணக்கோட்பாட்டை உருவாக்கியபோது,
முடிவுறாத ஒரு தொகுப்பு என்ற கருத்தை ஏற்கத்தக்கதாக்குவது
அவரது நோக்கங்களில் ஒன்றாக இருந்தது.
இடைக்காலச் சிந்தனையாளர்களின் சொல்வழக்கில், இது மெய்யாக அமைந்த
ஒரு முடிவிலி ஆகும். இதில் ஒரு முக்கியப் பகுதி,
வெவ்வேறு கணங்களின் \emph{அளவை} அவர் கையாண்ட விதமாகும்.
$a$, $b$, $c$ ஆகியவை அனைத்தும் வெவ்வேறானவை எனில்,
$\{a, b, c\}$ என்ற கணம் $\{a, b\}$ ஐ விடப்
\emph{பெரியது} என்பது உள்ளுணர்வுக்கு இயல்பாகத் தோன்றுகிறது.
ஆனால் முடிவுறாத கணங்களைப் பற்றிய நிலை என்ன?
அவை அனைத்தும் ஒன்றுக்கொன்று சம அளவுடையவையா?
அவ்வாறு இல்லை என்பது தெரியவருகிறது.

% SEGMENT OLP-0028-S02
இங்கு முதல் முக்கியமான கருத்து பட்டியலாக்கம் ஆகும்.
ஒவ்வொரு முடிவுறு கணத்தின் எல்லா !!{element}s[acc] பட்டியலிடுவதன்
மூலம் அந்தக் கணத்தைப் பட்டியலிடலாம்.
பட்டியலே முடிவுறாததாக இருக்க அனுமதித்தால்,
சில முடிவுறாத கணங்களின் எல்லா !!{element}s[acc]யும்
பட்டியலிடலாம். இத்தகைய கணங்கள் !!{enumerable} கணங்கள் எனப்படும்.
சில முடிவுறாத கணங்கள் !!{enumerable} தன்மையற்றவை என்பதே
காண்டோரின் வியப்பூட்டும் முடிவு.
இந்த அத்தியாயத்தின் முடிவில் அந்த முடிவை முழுமையாகப் புரிந்துகொள்வோம்.

\end{document}
